\documentclass[a4paper,12pt]{article}
\usepackage[utf8]{inputenc}
\usepackage[T1]{fontenc}
\usepackage[ngerman]{babel}
\usepackage{geometry}
\geometry{margin=2.5cm}
\usepackage{hyperref}

\begin{document}

\section*{Physics-Informed Neural Operator for Learning Partial Differential Equations}

\textbf{Autoren:} Zongyi Li, Hongkai Zheng, Nikola Kovachki, David Jin, Haoxuan Chen,
Burigede Liu, Kamyar Azizzadenesheli, Anima Anandkumar \\
\textbf{Journal:} ACM/IMS Journal of Data Science, Vol.~1, No.~3, Article~9 \\
\textbf{Jahr:} 2024 \\
\textbf{DOI:} \href{https://doi.org/10.1145/3648506}{10.1145/3648506}

\subsection*{Motivation und Problemstellung}

Maschinenlernen-Methoden zur Lösung partieller Differentialgleichungen (PDEs) lassen sich
grob in zwei Kategorien einteilen: solche, die eine Lösungsfunktion approximieren (z.\,B.
Physics-Informed Neural Networks, PINNs), und solche, die den Lösungsoperator lernen
(z.\,B. Fourier Neural Operators, FNO). Beide Ansätze haben charakteristische Schwächen:
PINNs leiden unter einer schwierigen Optimierungslandschaft, die bei multiskaligen oder
zeitabhängigen Systemen häufig zu Konvergenzproblemen führt. FNOs hingegen sind rein
datengetrieben und benötigen ausreichend hochwertige Trainingsdaten, deren Erzeugung
mittels numerischer Löser teuer oder in manchen Szenarien sogar unmöglich sein kann.

\subsection*{Methodik}

Die Arbeit schlägt den \textit{Physics-Informed Neural Operator} (PINO) vor, der beide
Paradigmen kombiniert: PINO nutzt sowohl Trainingsdaten als auch physikalische
Verlustfunktionen in Form von PDE-Residuen, um den Lösungsoperator einer parametrischen
PDE-Familie zu erlernen. Das Training gliedert sich in zwei Phasen:

\begin{enumerate}
    \item \textbf{Operator-Lernphase:} PINO lernt den Lösungsoperator über mehrere
    Instanzen der PDE-Familie unter Verwendung von Trainingsdaten \emph{und}
    physikalischen Nebenbedingungen. Dabei wird gezielt eine Kombination aus
    niedrigauflösenden Trainingsdaten mit hochauflösenden PDE-Residuen eingesetzt.
    Dies ermöglicht eine hochauflösende Rekonstruktion des wahren Operators, selbst
    wenn nur grobe Trainingsdaten vorliegen.

    \item \textbf{Instanzweises Finetuning:} Der vortrainierte Operator wird als
    Ansatzfunktion genutzt und für eine spezifische PDE-Instanz anhand der
    physikalischen Verlustfunktion weiter optimiert. Ein optionaler Ankerverlust
    ($\mathcal{L}_{\mathrm{op}}$) verhindert dabei, dass das Modell zu weit vom
    vortrainierten Zustand abweicht, was die Optimierung stabilisiert.
\end{enumerate}

Die Ableitungen für den PDE-Verlust werden effizient im Fourier-Raum berechnet, was
gegenüber punktweiser automatischer Differentiation (wie bei PINNs) erhebliche
Rechenvorteile bietet. Für nicht-periodische Randbedingungen wird die
\textit{Fourier-Continuation}-Technik eingesetzt, bei der die Eingabe durch Zero-Padding
in einen periodischen Bereich eingebettet wird.

\subsection*{Ergebnisse}

PINO wird auf mehreren Standard-Benchmarks evaluiert, darunter die Burgers-Gleichung,
die Darcy-Strömung und die Navier-Stokes-Gleichungen (inkl. des chaotischen
Kolmogorov-Flows). Die zentralen Ergebnisse lassen sich wie folgt zusammenfassen:

\begin{itemize}
    \item \textbf{Zero-Shot-Superresolution:} PINO zeigt keine signifikante
    Genauigkeitseinbuße, wenn es bei höheren Auflösungen als im Training evaluiert
    wird -- ein Vorteil gegenüber UNet-basierten Modellen, die ohne explizite
    Interpolation keine Superresolution leisten können.

    \item \textbf{Datensparsamkeit:} Mit physikalischen Verlustfunktionen kann PINO
    bei einfacheren PDEs (Burgers, Darcy) auch gänzlich ohne gelabelte Daten trainiert
    werden. Beim Burgers-Problem erzielt PINO 0,38\,\% relativen $L_2$-Fehler,
    verglichen mit 1,38\,\% für PI-DeepONet.

    \item \textbf{Chaotische Strömungen:} Bei Kolmogorov-Flows (Re\,=\,500) übertrifft
    PINO klassische PINN-Varianten um den Faktor 20 (kleinerer Fehler) bei gleichzeitig
    25-facher Beschleunigung.

    \item \textbf{Transferlernen:} Ein auf einem Reynolds-Zahl-Regime trainierter Operator
    lässt sich durch instanzweises Finetuning effizient auf andere Reynolds-Zahlen
    (100--500) übertragen.

    \item \textbf{Inverse Probleme:} Das PDE-Residuum schränkt die Lösung auf die
    physikalisch zulässige Mannigfaltigkeit ein, was die Qualität der Inversion
    gegenüber rein datengetriebenen Methoden deutlich verbessert. PINO ist dabei
    ca.\ 3000-mal schneller als eine MCMC-basierte Referenzlösung.
\end{itemize}

\subsection*{Bezug zur Masterarbeit}

Diese Arbeit ist in mehrfacher Hinsicht relevant für die vorliegende Masterarbeit, die
sich mit physikalisch-informiertem maschinellen Lernen zur Vorhersage von Strömungsfeldern
um sedimentierende Kristalle in magmatischen Systemen beschäftigt.

\paragraph{FNO als Architekturkandidat.}
Die Masterarbeit vergleicht UNet- und FNO-Architekturen als gleichwertige Hauptbeiträge.
PINO baut auf dem FNO-Rahmen auf und zeigt dessen universelle Approximationseigenschaften
sowie Diskretisierungskonvergenz -- beides sind theoretische Grundlagen, die die Wahl von
FNO als Vergleichsarchitektur in der Masterarbeit motivieren. Insbesondere die
Superresolution-Eigenschaften des FNO stehen in direktem Kontrast zu den Limitierungen
von UNets, die ohne Interpolation nicht auf ungesehene Gitterauflösungen generalisieren
können.

\paragraph{Physikalisch-informierte Verlustfunktionen.}
In der Masterarbeit werden bereits Kontinuitätsbedingungen ($\nabla \cdot \mathbf{v} = 0$)
als physikalische Verlustterme implementiert. Der PINO-Ansatz motiviert, diese
Physikterme systematischer zu nutzen: Die Stokes-Gleichungen, die das
Strömungsregime bei sehr kleinen Reynolds-Zahlen beschreiben, könnten analog zum
PINO-Framework als PDE-Residuum in die Verlustfunktion integriert werden. Dies würde
über die reine Divergenzfreiheitsbedingung hinausgehen und die Impulserhaltung direkt
berücksichtigen.

\paragraph{Residuallernansatz.}
Der in der Masterarbeit verwendete Residuallernansatz ($\mathbf{v}_{\mathrm{total}} =
\mathbf{v}_{\mathrm{Stokes}} + \Delta\mathbf{v}_{\mathrm{learned}}$) ist konzeptuell
verwandt mit dem PINO-Finetuning: Beide Methoden nutzen eine analytische oder
vortrainierte Anfangslösung als Ausgangspunkt und lernen Korrekturen dazu. PINO
formalisiert diesen Gedanken durch den Ankerverlust $\mathcal{L}_{\mathrm{op}}$, der
das Modell in der Nähe der Ausgangslösung hält.

\paragraph{Generalisierung auf neue Kristallkonfigurationen.}
Die zentrale Forschungsfrage der Masterarbeit -- ob ein auf $N$ Kristallen trainiertes
Modell auf $N{+}m$ Kristalle generalisiert -- ist direkt mit der in PINO untersuchten
Fähigkeit zur Generalisierung jenseits der Trainingsverteilung verwandt. PINOs Ansatz,
PDE-Residuen als Regularisierung zu nutzen, zeigt, dass physikalische Nebenbedingungen
die Generalisierung signifikant verbessern können. Dies legt nahe, dass eine stärkere
Gewichtung des Divergenz- oder Stokes-Residuums in der Masterarbeit die
Generalisierungsfähigkeit auf unbekannte Kristallanzahlen verbessern könnte.

\paragraph{Abgrenzung.}
PINO ist primär für zeitabhängige oder räumlich kontinuierliche PDE-Familien ausgelegt
und setzt periodische oder einfach strukturierte Randbedingungen voraus. Die in der
Masterarbeit betrachteten Strömungsfelder um Kristalle weisen komplexe,
geometrieabhängige Randbedingungen auf, die eine direkte Übertragung des PINO-Frameworks
erschweren. Dennoch liefert PINO wichtige methodische Impulse, insbesondere für zukünftige
Arbeiten, die über den datengetriebenen Vergleich von UNet und FNO hinausgehen.

\end{document}